% Lessons 38--49: intervals, graphs, linked lists, DP, windows, and bits.

\problemstart{Non-overlapping Intervals}{435}{Medium}{Greedy by earliest finish}{Intervals}

\subsection{The Problem}
Return the minimum number of intervals to remove so the remaining intervals
do not overlap.

\subsection{Keep the Interval That Finishes First}
Sort by end time. Whenever the next interval starts before the last kept end,
remove it. Choosing the earliest possible end leaves the greatest amount of
space for all future intervals.

\begin{invariantbox}{the kept intervals finish as early as possible}
After scanning a prefix, the algorithm has kept the maximum possible number
of compatible intervals from that prefix, and the final kept interval has the
smallest achievable end among choices of that size.
\end{invariantbox}

\begin{visualbox}{An early finish keeps more future choices}
\centering
\begin{tikzpicture}[font=\small,x=8mm]
  \node[anchor=east,text=teal,font=\bfseries] at (0.7,1.25) {keep};
  \draw[line width=4pt,teal] (1,1.25)--(3,1.25);
  \node[above=2mm] at (2,1.25) {\([1,3]\)};
  \draw[line width=4pt,navy] (3.25,1.25)--(5,1.25);
  \draw[line width=4pt,navy] (5.25,1.25)--(7,1.25);
  \node[below=2mm,text=navy] at (5.1,1.25) {two later intervals still fit};

  \node[anchor=east,text=orange,font=\bfseries] at (0.7,-0.35) {remove};
  \draw[line width=4pt,orange] (1,-0.35)--(6,-0.35);
  \node[above=2mm] at (3.5,-0.35) {\([1,6]\)};
  \node[below=2mm,text=orange] at (3.5,-0.35) {this long interval blocks both};
\end{tikzpicture}
\end{visualbox}

\subsection{Pseudocode}
\begin{lstlisting}[style=pseudocode]
sort intervals by end time
keptEnd = end of first interval; removals = 0
for each remaining interval:
    if interval.start < keptEnd: increment removals
    else: keptEnd = interval.end
return removals
\end{lstlisting}

\subsection{C++ Implementation}
\begin{lstlisting}[style=compactcode,caption={Complete C++ solution: Non-overlapping Intervals}]
#include <algorithm>
#include <vector>

using namespace std;

class Solution {
public:
    int eraseOverlapIntervals(vector<vector<int>>& intervals) {
        sort(intervals.begin(), intervals.end(),
             [](const vector<int>& a, const vector<int>& b) {
                 return a[1] < b[1];
             });
        int removals = 0;
        int lastEnd = intervals[0][1];
        for (int i = 1; i < static_cast<int>(intervals.size()); ++i) {
            if (intervals[i][0] < lastEnd) {
                ++removals;
            } else {
                lastEnd = intervals[i][1];
            }
        }
        return removals;
    }
};
\end{lstlisting}

\subsection{Why It Works and Complexity}
The standard exchange argument replaces any later-finishing kept interval with
the greedy one without harming compatibility. Sorting dominates at
\(O(n\log n)\); the scan uses \(O(1)\) extra state. Touching endpoints do
not overlap.

\lessonexpansion
  {For \texttt{[[1,2],[2,3],[3,4],[1,3]]}, keep \([1,2]\), then \([2,3]\),
   then \([3,4]\). The interval \([1,3]\) conflicts and is removed, producing
   the minimum removal count one.}
  {Among overlapping choices, keeping the interval with the earliest finish
   leaves at least as much room for every future interval as keeping a later
   finish. Repeating this exchange argument proves the greedy optimum.}
  {Check empty input, one interval, touching endpoints, identical intervals,
   nested intervals, and several conflicts sharing the same ending time.}
  {I sort by end time and keep each interval whose start is at least the last
   kept end. Every rejected interval overlaps the most future-friendly choice,
   so counting those rejections gives the minimum removals.}

\chaptercheckpoint
  {Remove as few intervals as possible to eliminate overlap.}
  {Intervals sorted by end and the end of the last kept interval.}
  {Sorting only by start and keeping a later-finishing conflict.}
  {$O(n\log n)$ time.}


\problemstart{Merge Intervals}{56}{Medium}{Sort and sweep}{Intervals}

\subsection{The Problem}
Merge every pair of overlapping intervals and return non-overlapping intervals.

\subsection{Build One Active Interval}
After sorting by start, compare the next start with the end of the last merged
interval. An overlap extends that end; otherwise the next interval begins a
new non-overlapping group.

\begin{invariantbox}{the output is merged and ends with the active group}
Before each comparison, all output intervals except possibly the last are
final and cannot overlap any later interval. The last interval equals the union of the current
overlapping group.
\end{invariantbox}

\begin{visualbox}{Overlapping ranges collapse into one range}
\centering
\begin{tikzpicture}[font=\small,>={Stealth[length=2.2mm,width=1.55mm]},x=8mm]
  \foreach \x in {0,...,7} \node[below] at (\x,0) {\scriptsize \x};
  \draw[->,navy] (0,0)--(7.4,0);
  \node[anchor=east,text=teal] at (0.7,1.6) {first};
  \draw[line width=5pt,teal] (1,1.6)--(3,1.6);
  \node[above=2mm] at (2,1.6) {$[1,3]$};
  \node[anchor=east,text=orange] at (0.7,.75) {next};
  \draw[line width=5pt,orange] (2,.75)--(6,.75);
  \node[above=2mm] at (4,.75) {$[2,6]$};
  \node[anchor=east,text=navy,font=\bfseries] at (0.7,-1.0) {output};
  \draw[line width=6pt,navy] (1,-1.0)--(6,-1.0);
  \node[below=2mm] at (3.5,-1.0) {$[1,6]$};
  \draw[densely dotted,ink!45] (1,-.75)--(1,1.35) (6,-.75)--(6,.5);
\end{tikzpicture}
\end{visualbox}

\subsection{Pseudocode}
\begin{lstlisting}[style=pseudocode]
sort intervals by start time
merged = empty
for interval in sorted order:
    if merged is empty or interval starts after merged.back ends:
        append interval
    else: extend merged.back end
return merged
\end{lstlisting}

\subsection{C++ Implementation}
\begin{lstlisting}[style=compactcode,caption={Complete C++ solution: Merge Intervals}]
#include <algorithm>
#include <vector>

using namespace std;

class Solution {
public:
    vector<vector<int>> merge(
            vector<vector<int>>& intervals) {
        sort(intervals.begin(), intervals.end());
        vector<vector<int>> merged;
        for (const auto& interval : intervals) {
            if (merged.empty() || merged.back()[1] < interval[0]) {
                merged.push_back(interval);
            } else {
                merged.back()[1] = max(merged.back()[1], interval[1]);
            }
        }
        return merged;
    }
};
\end{lstlisting}

\subsection{Time and Space}
Time is \(O(n\log n)\) for sorting and \(O(n)\) for the sweep. Excluding the
returned vector, sorting may use \(O(\log n)\) stack space. Equal endpoints
merge because the intervals overlap at that endpoint.

\lessonexpansion
  {Sort \texttt{[[1,3],[2,6],[8,10],[15,18]]}. Merge the first two into
   \([1,6]\); the next starts after 6, so append \([8,10]\), followed by
   \([15,18]\).}
  {Once intervals are sorted by start, only the final merged interval can
   overlap the next input. Extending that final end keeps the exact union;
   otherwise appending begins a new non-overlapping group.}
  {Check one interval, contained intervals, equal endpoints, repeated
   intervals, one chain of overlaps, and input with no overlaps.}
  {I sort by start time and compare each interval with the end of the current
   merged block. It either extends that block or begins a new one.}

\chaptercheckpoint
  {Return the union of overlapping intervals.}
  {Sorted intervals and the final merged interval.}
  {Forgetting to extend with the larger of the two ends.}
  {$O(n\log n)$ time.}


\problemstart{Insert Interval}{57}{Medium}{Three-phase interval scan}{Intervals}

\subsection{The Problem}
Insert one interval into a sorted, non-overlapping interval list and merge any
new overlaps.

\subsection{Before, Overlap, After}
Append intervals ending before the new interval. Merge every interval whose
start is at most the new end. Append the merged interval, then copy all
remaining intervals.

\begin{invariantbox}{processed intervals are final and sorted}
During the first phase, copied intervals lie strictly before the new range.
During the second, the mutable new interval equals the union of every overlap
seen so far. The remaining intervals begin after it once merging stops.
\end{invariantbox}

\begin{visualbox}{Insertion has three consecutive regions}
\centering
\begin{tikzpicture}[font=\small,x=8mm]
  \draw[->,navy] (0,0)--(10,0);
  \foreach \x in {0,2,4,6,8,10}{\node[below=1mm] at (\x,0) {\x};}
  \draw[line width=4pt,teal] (0.5,1.35)--(1.5,1.35);
  \node[above=2mm,text=teal] at (1,1.35) {copy before};
  \draw[line width=4pt,orange] (2,1.35)--(5,1.35);
  \draw[line width=4pt,orange,dashed] (4,0.75)--(8,0.75);
  \node[above=2mm,text=orange] at (5,1.35) {merge \([2,5]\) with \([4,8]\)};
  \draw[line width=4pt,navy] (8.5,1.35)--(9.5,1.35);
  \node[above=2mm,text=navy] at (9,1.35) {copy after};
  \draw[line width=5pt,teal] (2,-0.75)--(8,-0.75);
  \node[below=2mm,text=teal,font=\bfseries] at (5,-0.75) {insert merged interval \([2,8]\)};
\end{tikzpicture}
\end{visualbox}

\subsection{Pseudocode}
\begin{lstlisting}[style=pseudocode]
append intervals ending before newInterval starts
while intervals overlap newInterval:
    expand newInterval to their minimum start and maximum end
append the expanded newInterval
append every interval starting after it
return result
\end{lstlisting}

\subsection{C++ Implementation}
\begin{lstlisting}[style=compactcode,caption={Complete C++ solution: Insert Interval}]
#include <algorithm>
#include <vector>

using namespace std;

class Solution {
public:
    vector<vector<int>> insert(
            vector<vector<int>>& intervals,
            vector<int>& newInterval) {
        vector<vector<int>> result;
        int i = 0, n = static_cast<int>(intervals.size());
        while (i < n && intervals[i][1] < newInterval[0])
            result.push_back(intervals[i++]);
        while (i < n && intervals[i][0] <= newInterval[1]) {
            newInterval[0] = min(newInterval[0], intervals[i][0]);
            newInterval[1] = max(newInterval[1], intervals[i][1]);
            ++i;
        }
        result.push_back(newInterval);
        while (i < n) result.push_back(intervals[i++]);
        return result;
    }
};
\end{lstlisting}

\subsection{Time, Space, and Edge Cases}
Every interval is copied or merged once: \(O(n)\) time and \(O(n)\) output
space. Test insertion before all intervals, after all intervals, inside one
interval, and across several intervals.

\lessonexpansion
  {Insert \([4,8]\) into \([[1,2],[3,5],[6,7],[9,12]]\). Copy \([1,2]\).
   Merge \([3,5]\) to obtain \([3,8]\), then merge \([6,7]\) without changing
   the bounds. Append \([3,8]\), followed by \([9,12]\).}
  {Because the original intervals are sorted and do not overlap, they form three
   consecutive regions: strictly before, overlapping, and strictly after the
   inserted interval. Processing these regions in order produces a sorted,
   fully merged result.}
  {Test an empty list, insertion before or after everything, containment in
   either direction, touching endpoints, and one new interval spanning the
   complete list.}
  {I exploit the existing sorted order. I first copy intervals that end before
   the new one, then enlarge the new boundaries across every overlap, and
   finally copy the remaining suffix. Each interval is visited once.}

\chaptercheckpoint
  {Insert into an already sorted, non-overlapping interval list.}
  {A scan position and the expanding new interval.}
  {Using strict comparison when touching endpoints must merge.}
  {$O(n)$ time.}


\patternintoc{Grid and Graph Components}
\problemstart{Number of Islands}{200}{Medium}{Grid DFS}{Matrix and recursion}

\subsection{The Problem}
Count connected groups of land cells in a grid. Connections are horizontal or
vertical, not diagonal.

\subsection{Flood One Island at a Time}
When an unvisited land cell is found, increment the answer and flood-fill its
entire component. Marking visited cells as water prevents recounting and makes
the grid itself the visited structure.

\begin{invariantbox}{every discovered land component is erased exactly once}
After the scan passes a cell, no remaining land cell in its connected
component survives. So each future land discovery belongs to a new
island and increments the count exactly once.
\end{invariantbox}

\begin{visualbox}{DFS floods all four-directionally connected land}
\centering
\begin{tikzpicture}[font=\small]
  \node[arraycell,fill=softorange,draw=orange,very thick] (g00) at (0,0) {1};
  \node[arraycell,fill=softorange,draw=orange,very thick] (g01) at (1.05,0) {1};
  \node[arraycell,fill=white] (g02) at (2.10,0) {0};
  \node[arraycell,fill=white] (g03) at (3.15,0) {0};
  \node[arraycell,fill=softorange,draw=orange,very thick] (g10) at (0,-0.9) {1};
  \node[arraycell,fill=white] (g11) at (1.05,-0.9) {0};
  \node[arraycell,fill=white] (g12) at (2.10,-0.9) {0};
  \node[arraycell,fill=softteal,draw=teal,very thick] (g13) at (3.15,-0.9) {1};
  \node[arraycell,fill=white] (g20) at (0,-1.8) {0};
  \node[arraycell,fill=white] (g21) at (1.05,-1.8) {0};
  \node[arraycell,fill=softteal,draw=teal,very thick] (g22) at (2.10,-1.8) {1};
  \node[arraycell,fill=softteal,draw=teal,very thick] (g23) at (3.15,-1.8) {1};
  \node[above=4mm of g00,text=orange,font=\bfseries] {island 1};
  \node[right=7mm of g13,text=teal,font=\bfseries] {island 2};
\end{tikzpicture}
\end{visualbox}

\subsection{Pseudocode}
\begin{lstlisting}[style=pseudocode]
answer = 0
for every grid cell:
    if cell is land:
        increment answer
        DFS from it, changing all connected land to visited water
return answer
\end{lstlisting}

\subsection{C++ Implementation}
\begin{lstlisting}[style=compactcode,caption={Complete C++ solution: Number of Islands}]
#include <vector>

using namespace std;

class Solution {
    void flood(vector<vector<char>>& grid, int row, int col) {
        if (row < 0 || col < 0 || row >= static_cast<int>(grid.size()) ||
            col >= static_cast<int>(grid[0].size()) || grid[row][col] != '1')
            return;
        grid[row][col] = '0';
        flood(grid, row + 1, col); flood(grid, row - 1, col);
        flood(grid, row, col + 1); flood(grid, row, col - 1);
    }

public:
    int numIslands(vector<vector<char>>& grid) {
        int islands = 0;
        for (int r = 0; r < static_cast<int>(grid.size()); ++r) {
            for (int c = 0; c < static_cast<int>(grid[0].size()); ++c) {
                if (grid[r][c] == '1') {
                    ++islands;
                    flood(grid, r, c);
                }
            }
        }
        return islands;
    }
};
\end{lstlisting}

\subsection{Time, Space, and Safety}
Each cell is examined a constant number of times: \(O(mn)\) time. Recursion
can use \(O(mn)\) stack space for one large island. The rules provide a
non-empty grid; an iterative stack can avoid call-stack limits in production.

\lessonexpansion
  {In the grid \texttt{[[1,1,0],[1,0,0],[0,0,1]]}, the scan first discovers
   the upper-left land cell. DFS erases the three connected cells, so none can
   be counted again. The final lower-right land cell starts a second DFS and
   produces the answer two.}
  {A new island is counted only at an unvisited land cell. Flood fill reaches
   exactly its four-directional connected component and marks every member,
   so later scan positions cannot count that component again.}
  {Check a single water cell, a single land cell, all land, all water, a thin
   row, a thin column, diagonal-only contact, and a large snake-shaped island.}
  {I scan every cell and launch DFS only from unvisited land. That one DFS
   consumes the complete island. Since each cell changes from land to visited
   at most once, the total work is linear in the grid size.}

\chaptercheckpoint
  {Count connected land regions in a four-neighbor grid.}
  {The grid doubles as visited state during flood fill.}
  {Including diagonal cells or marking after recursion.}
  {$O(mn)$ time and up to $O(mn)$ stack space.}


\problemstart{Number of Connected Components in an Undirected Graph}{323}{Medium}{Graph DFS}{Adjacency list}

\subsection{The Problem}
Given \(n\) vertices and undirected edges, return the number of connected
components.

\subsection{Start One Traversal per Component}
Build an adjacency list. Each unvisited vertex starts a DFS that marks every
vertex reachable from it; this single start represents one component.

\begin{invariantbox}{visited vertices belong to already counted components}
After a DFS finishes, every vertex in that connected component is marked.
So the next unvisited vertex cannot belong to any component already counted.
\end{invariantbox}

\begin{visualbox}{Two DFS starts reveal two components}
\centering
\begin{tikzpicture}[font=\small]
  \node[trienode,visited] (a) {0}; \node[trienode,right=13mm of a] (b) {1};
  \node[trienode,below=11mm of b] (c) {2};
  \draw[-,thick,teal] (a)--(b)--(c);
  \node[trienode,active,right=30mm of b] (d) {3};
  \node[trienode,right=13mm of d] (e) {4};
  \draw[-,thick,orange] (d)--(e);
  \node[below=9mm of c]{component 1}; \node[below=20mm of d]{component 2};
\end{tikzpicture}
\end{visualbox}

\subsection{Pseudocode}
\begin{lstlisting}[style=pseudocode]
build an undirected adjacency list
visited = all false; components = 0
for each vertex:
    if vertex is unvisited:
        increment components
        DFS and mark its entire component
return components
\end{lstlisting}

\subsection{C++ Implementation}
\begin{lstlisting}[style=compactcode,caption={Complete C++ solution: Connected Components}]
#include <vector>

using namespace std;

class Solution {
    void visit(int node, const vector<vector<int>>& graph,
               vector<bool>& seen) {
        seen[node] = true;
        for (int neighbor : graph[node])
            if (!seen[neighbor]) visit(neighbor, graph, seen);
    }

public:
    int countComponents(int n, vector<vector<int>>& edges) {
        vector<vector<int>> graph(n);
        for (const auto& edge : edges) {
            graph[edge[0]].push_back(edge[1]);
            graph[edge[1]].push_back(edge[0]);
        }
        vector<bool> seen(n, false);
        int components = 0;
        for (int node = 0; node < n; ++node) {
            if (!seen[node]) {
                ++components;
                visit(node, graph, seen);
            }
        }
        return components;
    }
};
\end{lstlisting}

\subsection{Time and Space}
Building and traversing the graph costs \(O(V+E)\) time and space. Isolated
vertices are valid one-vertex components. Because the graph is undirected,
each edge must be added in both directions.

\lessonexpansion
  {With \(n=5\) and edges \texttt{[[0,1],[1,2],[3,4]]}, DFS from 0 marks
   \(0,1,2\) as one component. The next unvisited vertex is 3, whose DFS also
   marks 4. Exactly two DFS launches occur, so the answer is two.}
  {A DFS started at an unvisited vertex reaches exactly its connected
   component. Marking that entire set prevents recounting it, while every
   later unvisited vertex must belong to a different component.}
  {Check no edges, one vertex, a fully connected graph, duplicate edges,
   self-loops, and several isolated vertices.}
  {I build an undirected adjacency list and launch DFS from each still-unseen
   vertex. Each launch discovers one complete component, giving linear time
   in the vertices and edges.}

\chaptercheckpoint
  {Count disconnected regions of an undirected graph.}
  {An adjacency list and a visited flag per vertex.}
  {Adding each undirected edge in only one direction.}
  {$O(V+E)$ time and space.}


\problemstart{Longest Consecutive Sequence}{128}{Medium}{Hash-set starts}{Hash set}

\subsection{The Problem}
Return the length of the longest run of consecutive integer values in an
unsorted array. The algorithm should run in linear average time.

\subsection{Count Only from Sequence Starts}
Put every value in a set. A value begins a sequence only when its previous
is absent. From such starts, count upward until the sequence ends. Interior
values never launch redundant scans.

\begin{invariantbox}{each sequence is expanded exactly once}
Only the smallest value of a consecutive run lacks a previous. So
one scan accounts for the entire run, and every value participates in at most
one forward expansion.
\end{invariantbox}

\begin{visualbox}{The missing previous identifies a start}
\centering
\begin{tikzpicture}[font=\small,>={Stealth[length=2.2mm,width=1.55mm]}]
  \foreach \x/\v in {0/100,1/4,2/200,3/1,4/3,5/2}{
    \node[arraycell] (v\x) at (1.05*\x,1.0) {\v};
  }
  \foreach \x/\v in {0/1,1/2,2/3,3/4}{
    \node[arraycell,fill=softteal] (s\x) at (1.05*\x,0) {\v};
  }
  \node[left=9mm of s0,text=orange]{0 absent};
  \draw[flowarrow,orange] ([xshift=-7mm]s0.west)--(s0);
  \node[right=8mm of s3]{length 4};
\end{tikzpicture}
\end{visualbox}

\subsection{Pseudocode}
\begin{lstlisting}[style=pseudocode]
values = hash set of all numbers
answer = 0
for value in values:
    if value - 1 is absent:
        length = 1
        while value + length exists: increment length
        answer = max(answer, length)
return answer
\end{lstlisting}

\subsection{C++ Implementation}
\begin{lstlisting}[style=compactcode,caption={Complete C++ solution: Longest Consecutive Sequence}]
#include <algorithm>
#include <unordered_set>
#include <vector>

using namespace std;

class Solution {
public:
    int longestConsecutive(vector<int>& nums) {
        unordered_set<int> values(nums.begin(), nums.end());
        int answer = 0;
        for (int value : values) {
            if (!values.count(value - 1)) {
                int current = value, length = 1;
                while (values.count(current + 1)) {
                    ++current;
                    ++length;
                }
                answer = max(answer, length);
            }
        }
        return answer;
    }
};
\end{lstlisting}

\subsection{Time and Space}
Set operations are \(O(1)\) average, and every distinct value is advanced
through once, giving \(O(n)\) average time and \(O(n)\) space. Duplicates are
removed automatically by the set.

\lessonexpansion
  {For \texttt{[100,4,200,1,3,2]}, only 100, 200, and 1 lack predecessors.
   The scan starting at 1 advances through 2, 3, and 4, producing length four;
   values 2, 3, and 4 never start duplicate scans.}
  {Every consecutive sequence has exactly one value with no previous.
   Starting only there counts every member once and prevents the quadratic
   repeated work caused by expanding from interior values.}
  {Check empty input, duplicates, negative values, one long sequence, several
   tied sequences, and values near integer limits.}
  {I place values in a hash set and expand only from values whose previous
   is absent. Those are exactly the starts of maximal sequences, so the
   total expected work is linear.}

\chaptercheckpoint
  {Find consecutive values without sorting.}
  {A set and values whose previous is absent.}
  {Starting a forward scan from every interior value.}
  {$O(n)$ average time and $O(n)$ space.}


\patternintoc{Linked-List Merging}
\problemstart{Merge Two Sorted Lists}{21}{Easy}{Two-pointer merge}{Linked lists}

\subsection{The Problem}
Merge two sorted singly linked lists by reusing their nodes and return the
head of the sorted result.

\subsection{Attach the Smaller Front Node}
A dummy node removes special handling for the first attachment. At each step,
link the smaller current node after the tail and advance that list. Once one
list ends, append the other remainder.

\begin{invariantbox}{the merged prefix is sorted and complete}
The nodes before \texttt{tail} are exactly the smallest consumed nodes from
both lists in sorted order. The two current pointers identify the smallest
remaining choice from each input.
\end{invariantbox}

\begin{visualbox}{Two sorted fronts feed one output tail}
\centering
\begin{tikzpicture}[font=\small,>={Stealth[length=2.2mm,width=1.55mm]}]
  \node[dsnode,active] (a1) at (0,1.3) {1};
  \node[dsnode] (a2) at (2.0,1.3) {2};
  \node[dsnode] (a3) at (4.0,1.3) {4};
  \node[left=6mm of a1,text=navy] {list 1};
  \draw[flowarrow] (a1)--(a2)--(a3);
  \node[dsnode,active] (b1) at (0,0) {1};
  \node[dsnode] (b2) at (2.0,0) {3};
  \node[dsnode] (b3) at (4.0,0) {4};
  \node[left=6mm of b1,text=navy] {list 2};
  \draw[flowarrow] (b1)--(b2)--(b3);
  \node[left=18mm of b1,text=orange,font=\bfseries,align=right]
    {compare the\\two front nodes};
  \node[dsnode] (dummy) at (-2.0,-1.55) {dummy};
  \node[dsnode,visited] (out1) at (0,-1.55) {1};
  \draw[flowarrow,teal] (dummy)--(out1);
  \draw[flowarrow,orange] (b1.south)--node[edgelabel,right]{attach} (out1.north);
  \node[right=7mm of out1,text=teal,font=\bfseries] {tail};
\end{tikzpicture}
\end{visualbox}

\subsection{Pseudocode}
\begin{lstlisting}[style=pseudocode]
dummy tail starts an empty result
while both lists are non-empty:
    attach the smaller front node; advance that list
attach the remaining non-empty list
return dummy.next
\end{lstlisting}

\subsection{C++ Implementation}
\begin{lstlisting}[style=compactcode,caption={Complete C++ solution: Merge Two Sorted Lists}]
using namespace std;

class Solution {
public:
    ListNode* mergeTwoLists(ListNode* list1, ListNode* list2) {
        ListNode dummy(0);
        ListNode* tail = &dummy;
        while (list1 != nullptr && list2 != nullptr) {
            if (list1->val <= list2->val) {
                tail->next = list1;
                list1 = list1->next;
            } else {
                tail->next = list2;
                list2 = list2->next;
            }
            tail = tail->next;
        }
        tail->next = (list1 != nullptr) ? list1 : list2;
        return dummy.next;
    }
};
\end{lstlisting}

\subsection{Time and Space}
Each node is linked once: \(O(m+n)\) time and \(O(1)\) extra space. The
dummy node is stack allocated and is not part of the returned list.

\lessonexpansion
  {Merge \texttt{1->2->4} with \texttt{1->3->4}. Attach 1 from the first
   list, 1 from the second, then 2, 3, and 4. When one list becomes empty,
   attach the remaining suffix directly.}
  {At every step the smaller front node is the smallest value not yet emitted.
   Attaching it keeps sorted order, and once one list is empty the other
   suffix is already sorted and no longer competes with another value.}
  {Check two null lists, one null list, equal front values, one list entirely
   smaller, and lists containing negative or repeated values.}
  {I use a dummy head and one tail pointer. The smaller current node is linked
   after the tail, its source advances, and the final remaining suffix is
   attached in constant time.}

\chaptercheckpoint
  {Combine two already sorted linked sequences.}
  {Two current nodes and one output tail.}
  {Advancing the output tail without advancing the chosen input.}
  {$O(m+n)$ time and $O(1)$ space.}


\problemstart{Merge K Sorted Lists}{23}{Hard}{Minimum-node heap}{Linked lists and min heap}

\subsection{The Problem}
Merge \(k\) sorted linked lists into one sorted list.

\subsection{Expose One Choice per List}
The smallest remaining node must be one of the current list heads. Store those
heads in a min heap, remove the smallest, attach it, and push its next.

\begin{invariantbox}{the heap contains every possible next output node}
For each non-empty unmerged list suffix, the heap contains its first node.
Because each suffix is sorted, the minimum heap entry is globally the smallest
remaining node and is safe to append.
\end{invariantbox}

\begin{visualbox}{One heap entry represents each list frontier}
\centering
\begin{tikzpicture}[font=\small,>={Stealth[length=2.2mm,width=1.55mm]}]
  \node[trienode,active] (h) {1};
  \node[trienode,below left=11mm and 9mm of h] (a) {1};
  \node[trienode,below right=11mm and 9mm of h] (b) {2};
  \draw[-,thick,navy] (h)--(a); \draw[-,thick,navy] (h)--(b);
  \node[left=23mm of h,dsnode] (lists) {list heads\\1, 1, 2};
  \draw[flowarrow] (lists)--(h);
  \node[right=20mm of h,dsnode,visited] (out) {append 1\\push next};
  \draw[flowarrow,teal] (h)--(out);
\end{tikzpicture}
\end{visualbox}

\subsection{Pseudocode}
\begin{lstlisting}[style=pseudocode]
push every non-null list head into a min heap
dummy tail starts the result
while heap is not empty:
    node = pop smallest
    attach node to result
    if node.next exists: push node.next
return dummy.next
\end{lstlisting}

\subsection{C++ Implementation}
\begin{lstlisting}[style=compactcode,caption={Complete C++ solution: Merge K Sorted Lists}]
#include <queue>
#include <vector>

using namespace std;

class Solution {
    struct Compare {
        bool operator()(ListNode* a, ListNode* b) const {
            return a->val > b->val;
        }
    };
public:
    ListNode* mergeKLists(vector<ListNode*>& lists) {
        priority_queue<ListNode*, vector<ListNode*>, Compare> heap;
        for (ListNode* head : lists) if (head != nullptr) heap.push(head);
        ListNode dummy(0);
        ListNode* tail = &dummy;
        while (!heap.empty()) {
            ListNode* node = heap.top();
            heap.pop();
            tail->next = node;
            tail = node;
            if (node->next != nullptr) heap.push(node->next);
        }
        return dummy.next;
    }
};
\end{lstlisting}

\subsection{Time and Space}
For \(N\) total nodes, each heap operation costs \(O(\log k)\), producing
\(O(N\log k)\) time and \(O(k)\) heap space. Null list heads are never pushed.

\lessonexpansion
  {For lists \texttt{1->4->5}, \texttt{1->3->4}, and \texttt{2->6}, the heap
   begins with \(1,1,2\). Pop the first 1 and push 4; pop the second 1 and push
   3; then pop 2 and push 6. Continuing yields
   \texttt{1->1->2->3->4->4->5->6}.}
  {Every remaining list is sorted, so its head is its smallest unmerged node.
   The smallest remaining node must be among the heap entries. Popping
   it is safe, and pushing its next restores one frontier per list.}
  {Check zero lists, all-null lists, one list, duplicate values, lists of very
   different lengths, and negative values.}
  {I place one current node from each list into a min heap. The heap exposes
   the next global output value, after which I insert only that node's
   next. This uses \(O(k)\) extra space rather than storing all nodes.}

\chaptercheckpoint
  {Merge several sorted linked lists.}
  {A min heap containing one frontier node per non-empty list.}
  {Pushing every node initially and wasting linear heap space.}
  {$O(N\log k)$ time and $O(k)$ space.}


\patternintoc{String DP and Windows}
\problemstart{Longest Common Subsequence}{1143}{Medium}{Two-dimensional DP}{Strings and DP table}

\subsection{The Problem}
Return the length of the longest sequence appearing in both strings in the
same relative order. Characters need not be consecutive.

\subsection{Prefix Update Rule}
Let \(dp[i][j]\) be the LCS length of the first \(i\) characters of the first
string and the first \(j\) of the second. Equal final characters extend the
diagonal state; otherwise skip one final character from either string.

\begin{invariantbox}{each cell solves one pair of prefixes}
When \(dp[i][j]\) is computed, the top, left, and diagonal previous cells
already contain correct answers for smaller prefixes. The update rule covers
whether the two final characters are paired or one is excluded.
\end{invariantbox}

\begin{visualbox}{A match reads the diagonal DP cell}
\centering
\begin{tikzpicture}[font=\small,>={Stealth[length=2.2mm,width=1.55mm]}]
  \matrix[matrix of nodes,nodes={arraycell,minimum width=9mm},row sep=-\pgflinewidth,column sep=-\pgflinewidth] (m) {
    0 & 0 & 0 & 0 \\
    0 & 1 & 1 & 1 \\
    0 & 1 & 1 & 2 \\
    0 & 1 & 2 & 2 \\
  };
  \node[above=3mm of m-1-2]{a}; \node[above=3mm of m-1-3]{c}; \node[above=3mm of m-1-4]{e};
  \node[left=3mm of m-2-1]{a}; \node[left=3mm of m-3-1]{b}; \node[left=3mm of m-4-1]{c};
  % For the matching c/c cell, the diagonal predecessor is dp[2][1] = 1.
  \node[draw=orange,line width=1pt,fit=(m-3-2),inner sep=0.6pt] {};
  \node[draw=teal,line width=1pt,fit=(m-4-3),inner sep=0.6pt] {};
  \draw[->,very thick,orange] ([yshift=-1mm]m-3-2.south)
    to[out=-70,in=-110] node[edgelabel,below]{add 1 for matching \texttt{c}}
    ([yshift=-1mm]m-4-3.south);
  \node[dsnode,fill=softorange,right=16mm of m,yshift=7mm]
    {diagonal cell is \(1\)};
  \node[dsnode,fill=softteal,right=16mm of m,yshift=-8mm]
    {new cell is \(1+1=2\)};
\end{tikzpicture}
\end{visualbox}

\subsection{Pseudocode}
\begin{lstlisting}[style=pseudocode]
dp table starts at zero
for i from 1 to length(text1):
    for j from 1 to length(text2):
        if characters match: dp[i][j] = 1 + dp[i-1][j-1]
        else: dp[i][j] = max(dp[i-1][j], dp[i][j-1])
return dp[last row][last column]
\end{lstlisting}

\subsection{C++ Implementation}
\begin{lstlisting}[style=compactcode,caption={Complete C++ solution: Longest Common Subsequence}]
#include <algorithm>
#include <string>
#include <vector>

using namespace std;

class Solution {
public:
    int longestCommonSubsequence(string a, string b) {
        vector<vector<int>> dp(
            a.size() + 1, vector<int>(b.size() + 1, 0));
        for (int i = 1; i <= static_cast<int>(a.size()); ++i) {
            for (int j = 1; j <= static_cast<int>(b.size()); ++j) {
                if (a[i - 1] == b[j - 1]) dp[i][j] = 1 + dp[i - 1][j - 1];
                else dp[i][j] = max(dp[i - 1][j], dp[i][j - 1]);
            }
        }
        return dp[a.size()][b.size()];
    }
};
\end{lstlisting}

\subsection{Time and Space}
For lengths \(m,n\), time and table space are \(O(mn)\). Space can be reduced
to \(O(n)\) with two rows, but the full table makes dependencies and sequence
reconstruction clearer.

\lessonexpansion
  {For \texttt{abcde} and \texttt{ace}, the first match writes one at the
   \texttt{a}/\texttt{a} cell. Mismatches copy the maximum from above or left.
   Matches at \texttt{c} and \texttt{e} read the diagonal and add one, so the
   bottom-right answer is three.}
  {If the final characters match, a best subsequence can include that
   character and extend the solution for both shorter prefixes. If they differ,
   at least one final character is excluded, so the maximum of the top and
   left states is complete.}
  {Check one empty string, identical strings, no common characters, repeated
   characters, reversed order, and one string contained as a subsequence.}
  {I define each DP cell as the LCS length for two prefixes. Equal final
   characters use diagonal plus one; unequal characters choose the better
   result after skipping one side.}

\chaptercheckpoint
  {Characters keep order but may skip positions.}
  {A DP cell for every pair of string prefixes.}
  {Treating the problem as a consecutive substring.}
  {$O(mn)$ time and space.}


\problemstart{Minimum Window Substring}{76}{Hard}{Variable sliding window}{Character counts}

\subsection{The Problem}
Return the shortest substring of \texttt{s} containing every character of
\texttt{t}, including duplicate requirements. Return an empty string if none
exists.

\subsection{Expand Until Valid, Then Contract}
A requirement table stores how many of each character are still needed.
\texttt{missing} counts the total required occurrences not yet satisfied.
Once it reaches zero, move the left boundary while the valid condition remains and record
the shortest window.

\begin{invariantbox}{missing counts unsatisfied required occurrences}
After adding or removing a boundary character, the requirement count and
\texttt{missing} exactly describe whether the active window covers the target
multiset. The contraction loop ends at the smallest valid window for that
right endpoint.
\end{invariantbox}

\begin{visualbox}{The window contracts around the required multiset}
\centering
\begin{tikzpicture}[font=\small,>={Stealth[length=2.2mm,width=1.55mm]}]
  \foreach \x/\v in {0/A,1/D,2/O,3/B,4/E,5/C,6/O,7/D,8/E,9/B,10/A,11/N,12/C}{
    \node[arraycell,minimum width=8mm] (w\x) at (0.78*\x,0) {\v};
  }
  \node[draw=orange,very thick,fit=(w9)(w10)(w11)(w12),inner sep=1mm] {};
  \node[below=8mm of w10]{\texttt{BANC} covers \texttt{ABC}};
  \node[above=8mm of w0,text=navy] (oldleft) {old left edge};
  \node[above=8mm of w9,text=teal,font=\bfseries] (newleft) {new left edge};
  \draw[flowarrow] (oldleft.south)--(w0.north);
  \draw[flowarrow,teal] (newleft.south)--(w9.north);
  \draw[->,teal,very thick] (oldleft.east)--node[edgelabel,above]{move right} (newleft.west);
\end{tikzpicture}
\end{visualbox}

\subsection{Pseudocode}
\begin{lstlisting}[style=pseudocode]
count required characters and number of satisfied kinds
for each right boundary:
    add its character to the window
    while all required kinds are satisfied:
        record the window if it is shorter
        remove the left character and advance left
return the best recorded substring
\end{lstlisting}

\subsection{C++ Implementation}
\begin{lstlisting}[style=compactcode,caption={Complete C++ solution: Minimum Window Substring}]
#include <array>
#include <limits>
#include <string>

using namespace std;

class Solution {
public:
    string minWindow(string s, string t) {
        array<int, 256> need{};
        for (unsigned char ch : t) ++need[ch];
        int missing = static_cast<int>(t.size());
        int left = 0, bestStart = 0, bestLength = numeric_limits<int>::max();
        for (int right = 0; right < static_cast<int>(s.size()); ++right) {
            unsigned char in = static_cast<unsigned char>(s[right]);
            if (need[in] > 0) --missing;
            --need[in];
            while (missing == 0) {
                if (right - left + 1 < bestLength) {
                    bestLength = right - left + 1;
                    bestStart = left;
                }
                unsigned char out = static_cast<unsigned char>(s[left++]);
                ++need[out];
                if (need[out] > 0) ++missing;
            }
        }
        return bestLength == numeric_limits<int>::max()
             ? "" : s.substr(bestStart, bestLength);
    }
};
\end{lstlisting}

\subsection{Time and Space}
Each boundary advances at most \(n\) times, so time is \(O(n+|t|)\) and the
fixed character table is \(O(1)\). Duplicates in \texttt{t} are handled by
counts rather than a set.

\lessonexpansion
  {For \texttt{s=ADOBECODEBANC} and \texttt{t=ABC}, expand until the first
   valid window \texttt{ADOBEC}. Contracting loses the required A, so expansion
   resumes. Near the end, the valid condition is restored and repeated contraction
   produces \texttt{BANC}, the shortest length-four window.}
  {The requirement array and \texttt{missing} exactly encode whether the
   current window covers the target multiset. For each right endpoint, the
   inner loop records every valid left contraction and stops immediately after
   removing a required occurrence. So the shortest valid window is tested.}
  {Check an empty target if the interface permits it, a target longer than the
   source, repeated required characters such as \texttt{AABC}, exact equality,
   no solution, and a one-character answer.}
  {I expand a counted window until all required occurrences are present, then
   contract it as far as the valid condition allows. Counts are essential because a set
   cannot represent duplicate requirements. Both boundaries move only forward.}

\chaptercheckpoint
  {Find the smallest window covering a character multiset.}
  {Two boundaries, remaining counts, and total missing occurrences.}
  {Using distinct-character counts and losing duplicate requirements.}
  {$O(n+|t|)$ time and $O(1)$ alphabet space.}


\patternintoc{Bit Manipulation}
\problemstart{Missing Number}{268}{Easy}{XOR cancellation}{Array and bits}

\subsection{The Problem}
An array contains \(n\) distinct values chosen from \([0,n]\). Return the one
missing value.

\subsection{Cancel Equal Values with XOR}
XOR every expected value \(0\) through \(n\) and every array value. Each
present value occurs twice and cancels to zero; only the missing value remains.

\begin{invariantbox}{the running value is the XOR imbalance so far}
After processing index \(i\), the running value equals the XOR of the expected
indexes processed and the array values processed, together with \(n\). Values
appearing in both groups cancel regardless of order.
\end{invariantbox}

\begin{visualbox}{Pairs cancel and expose the missing value}
\centering
\begin{tikzpicture}[font=\small]
  \node[anchor=east,font=\bfseries,text=navy] at (-0.75,0.65) {expected};
  \node[anchor=east,font=\bfseries,text=teal] at (-0.75,-0.65) {in array};
  \foreach \x/\v in {0/0,1/1,2/2,3/3}{
    \node[arraycell,minimum width=11mm] (e\x) at (1.25*\x,0.65) {\v};
  }
  \node[arraycell,minimum width=11mm] (a0) at (0,-0.65) {0};
  \node[arraycell,minimum width=11mm] (a1) at (1.25,-0.65) {1};
  \node[draw=orange,dashed,minimum width=11mm,minimum height=8mm] at (2.5,-0.65) {?};
  \node[arraycell,minimum width=11mm] (a3) at (3.75,-0.65) {3};
  \foreach \x in {0,1,3}{
    \draw[teal,line width=1.1pt]
      ([xshift=-4mm,yshift=-3mm]e\x.center)--([xshift=4mm,yshift=3mm]e\x.center);
    \draw[teal,line width=1.1pt]
      ([xshift=-4mm,yshift=-3mm]a\x.center)--([xshift=4mm,yshift=3mm]a\x.center);
  }
  \node[dsnode,active,right=14mm of e3,yshift=-5mm]
    {only \(2\) has no pair\\missing number \(=2\)};
\end{tikzpicture}
\end{visualbox}

\subsection{Pseudocode}
\begin{lstlisting}[style=pseudocode]
answer = n
for i from 0 to n-1:
    answer = answer XOR i XOR nums[i]
return answer
\end{lstlisting}

\subsection{C++ Implementation}
\begin{lstlisting}[style=compactcode,caption={Complete C++ solution: Missing Number}]
#include <vector>

using namespace std;

class Solution {
public:
    int missingNumber(vector<int>& nums) {
        int answer = static_cast<int>(nums.size());
        for (int i = 0; i < static_cast<int>(nums.size()); ++i)
            answer ^= i ^ nums[i];
        return answer;
    }
};
\end{lstlisting}

\subsection{Time and Space}
Time is \(O(n)\), space is \(O(1)\), and XOR avoids arithmetic overflow. The
initial \(n\) is required because the loop indexes provide only \(0\) through
\(n-1\).

\lessonexpansion
  {For \texttt{[3,0,1]}, begin with \(answer=3\). XOR the pairs
   \((0,3),(1,0),(2,1)\). Equal values cancel, leaving 2, the only number that
   appears in the expected range but not in the array.}
  {XOR is associative and self-canceling. Combining every expected value
   \(0\ldots n\) with every actual value removes all matched pairs and leaves
   exactly the missing value.}
  {Check a missing zero, a missing \(n\), one-element input, shuffled input,
   and the largest permitted range.}
  {I XOR all indexes, all array values, and the endpoint \(n\). Every present
   number cancels with its expected copy, leaving the missing number with
   constant extra space.}

\chaptercheckpoint
  {One distinct value is absent from a complete integer range.}
  {An XOR running value over expected and actual values.}
  {Forgetting to include the endpoint $n$.}
  {$O(n)$ time and $O(1)$ space.}


\problemstart{Number of 1 Bits}{191}{Easy}{Clear lowest set bit}{Unsigned integer}

\subsection{The Problem}
Return the number of set bits in the binary representation of a 32-bit
unsigned integer.

\subsection{Remove One Set Bit per Iteration}
The expression \(n\mathbin{\&}(n-1)\) clears the lowest set bit. Repeating it
until zero makes the number of iterations equal to the number of ones.

\begin{invariantbox}{each loop removes exactly one one-bit}
Subtracting one flips the lowest one to zero and every lower zero to one.
ANDing with the original clears that one while preserving every higher bit.
\end{invariantbox}

\begin{visualbox}{Brian Kernighan's bit-clearing step}
\centering
\begin{tikzpicture}[font=\ttfamily\small]
  \matrix[matrix of nodes,nodes={minimum width=6mm,minimum height=7mm,
    draw=navy,fill=softblue},column sep=-\pgflinewidth,row sep=3mm] (bits) {
    1&0&1&1&0&0&0&0\\
    1&0&1&0&1&1&1&1\\
    1&0&1&0&0&0&0&0\\};
  \node[left=5mm of bits-1-1] {$n$};
  \node[left=5mm of bits-2-1] {$n-1$};
  \node[left=5mm of bits-3-1,font=\sffamily\bfseries] {AND};
  \node[draw=orange,very thick,fit=(bits-1-4)(bits-3-4),inner sep=1mm] {};
  \node[right=7mm of bits-3-8,font=\sffamily,align=left,text=orange]
    {lowest 1\\is cleared};
\end{tikzpicture}
\end{visualbox}

\subsection{Pseudocode}
\begin{lstlisting}[style=pseudocode]
count = 0
while n is not zero:
    n = n AND (n - 1)   // clear its lowest set bit
    increment count
return count
\end{lstlisting}

\subsection{C++ Implementation}
\begin{lstlisting}[style=compactcode,caption={Complete C++ solution: Number of 1 Bits}]
#include <cstdint>

using namespace std;

class Solution {
public:
    int hammingWeight(uint32_t n) {
        int count = 0;
        while (n != 0) {
            n &= n - 1;
            ++count;
        }
        return count;
    }
};
\end{lstlisting}

\subsection{Time and Space}
The loop runs once per set bit, at most 32 times: \(O(1)\) for fixed-width
input and \(O(1)\) space. Using an unsigned type avoids sign-extension issues.

\lessonexpansion
  {For binary \texttt{1011000}, subtracting one gives \texttt{1010111}; AND
   clears the lowest set bit to obtain \texttt{1010000}. Repeating produces
   \texttt{1000000} and then zero, so the count is three.}
  {For any nonzero unsigned integer, \(n\mathbin{\&}(n-1)\) removes exactly
   its lowest set bit and changes no higher set bit. The number of iterations
   equals the number of set bits.}
  {Check zero, one, all 32 bits set, only the highest bit set, and alternating
   bit patterns.}
  {I repeatedly clear the lowest set bit with \(n\mathbin{\&}(n-1)\) and count
   how many times that operation is possible.}

\chaptercheckpoint
  {Count set bits in a fixed-width integer.}
  {The current bit pattern and a counter.}
  {Using signed right shift for values with the top bit set.}
  {At most 32 iterations and $O(1)$ space.}